\documentclass{beamer}

\usetheme{Antibes}

% Loading preambulo
% PREAMBULO ====================================================================

% Packages ---------------------------------------------------------------------

% Package to Portuguese language
\usepackage[brazilian]{babel}
% Packages to Figures
\usepackage{graphicx}
\usepackage{tikz}
% Packages to math symbols and expressions
\usepackage{amsfonts, amssymb, amsmath}
% Packages to insert code
\usepackage{listings} 
\usepackage{verbatim}
% Package to justify text
\usepackage[document]{ragged2e}
% Package to manage the bibliography
\usepackage[backend=biber, style=numeric, sorting=none]{biblatex}
% Package to facilitate quotations
\usepackage{csquotes}
% Package to use multicols
\usepackage{multicol}
% Packages for url
%	The package libertinus already loads hyperref
%\usepackage[colorlinks=true]{hyperref}
% Font
%   Fira
%\usepackage[sfdefault, lf]{FiraSans}
%   Libertinus
\usepackage{libertinus}
\renewcommand*\familydefault{\sfdefault} %% Only if the base font of the document is to be sans serif
% Colors
\usepackage{xcolor}
\usepackage{colortbl}
% Table
\usepackage{tabularray}
% Fill line
\usepackage{xhfill}

% Configurations ---------------------------------------------------------------

\AtBeginSection{
    \begin{frame}{\secname}
        \tableofcontents[currentsection,hideallsubsections]
    \end{frame}
}

\AtBeginSubsection{
    \begin{frame}{\subsecname}
        \tableofcontents[currentsection,subsectionstyle=show/shaded/hide, subsubsectionstyle=hide]
    \end{frame}
}

\AtBeginSubsubsection{
    \begin{frame}{\subsubsecname}
        \tableofcontents[currentsection,subsectionstyle=show/shaded/hide,subsubsectionstyle=show/shaded/hide/hide]
    \end{frame}
}

% Numbering slides
\setbeamertemplate{footline}[frame number]{}

% Getting rid of bottom navigation bars
\setbeamertemplate{navigation symbols}{}

% Margins
\setbeamersize{text margin left=20pt, text margin right=20pt}

% Colors
\definecolor{burgundy}{RGB}{133,24,42}
\definecolor{bloodred}{RGB}{147,5,5}

\setbeamercolor{structure}{fg=burgundy}

\setbeamercolor*{palette primary}{use=structure,fg=white,bg=structure.fg}
%\setbeamercolor*{palette secondary}{use=structure,fg=white,bg=structure.fg!75!black}
%\setbeamercolor*{palette tertiary}{use=structure,fg=white,bg=structure.fg!50!black}
\setbeamercolor*{palette quaternary}{fg=white,bg=bloodred}

%	hyperref
\hypersetup{
	colorlinks=true,
	urlcolor=blue,
	breaklinks=true}

% New commands -----------------------------------------------------------------

\NewTblrTableCommand \aula{\SetCell{bg=green!60,fg=white}}
\NewTblrTableCommand \prova{\SetCell{bg=red!80,fg=white}}
\NewTblrTableCommand \feriado{\SetCell{bg=blue!50,fg=white}}

\newcommand{\esta}[1]{\textbf{\underline{#1}}}

% Title
\title[GP 05]{Gerência de Projetos}
% Subtitle
\subtitle{Revisão AV1}
% Author of the presentation
\author[E.J.R. Silva]{Evandro J.R. Silva}
% date of the presentation
\date{}

\begin{document}
	
\begin{frame}
	\titlepage
\end{frame}

\begin{frame}{Sumário}
	\tableofcontents
\end{frame}

%===============================================================================
% SECTION 1 ====================================================================
%===============================================================================
\section{Domínios de Desempenho de Projetos}

%-------------------------------------------------------------------------------
% SUBSEÇÃO 2.1 - DOMÍNIO DE DESEMPENHO DAS PARTES INTERESSADAS -----------------
%-------------------------------------------------------------------------------
\subsection{Partes Interessadas}

\begin{frame}{Partes Interessadas}
	\begin{enumerate}
		\justifying
		\item O processo de compreender as partes interessadas envolve \rule{50pt}{1pt}, \rule{50pt}{1pt}, \rule{50pt}{1pt} e priorização para determinar como engajá-las de forma eficaz.
		\begin{enumerate}[a]
			\item Identificação / Análise / Compreensão
			\item \alert<2>{Identificação / Análise / Compreensão}
			\item Registro / Monitoramento / Relatório
			\item Poder / Interesse / Influência
			\item Comunicação / Feedback / Reunião
		\end{enumerate}
	\end{enumerate}
\end{frame}

\begin{frame}{Partes Interessadas}
	\begin{enumerate}
		\setcounter{enumi}{1}
		\justifying
		\item O sucesso neste domínio é evidenciado quando as partes interessadas estão satisfeitas com o projeto e \rule{60pt}{1pt} são minimizados através do engajamento proativo.
		\begin{enumerate}[a]
			\item Mudanças
			\item \alert<2>{Conflitos}
			\item Riscos
			\item Custos
			\item Atrasos
		\end{enumerate}
	\end{enumerate}
\end{frame}

\begin{frame}{Partes Interessadas}
	\begin{enumerate}
		\setcounter{enumi}{2}
		\justifying
		\item Relacione a atividade de engajamento à sua descrição:
		\begin{columns}
			\begin{column}{0.28\textwidth}
				\begin{enumerate}[A]
					\item Monitorar
					\item Engajar
					\item Analisar
					\item Identificar
				\end{enumerate}
			\end{column}
			\begin{column}{0.72\textwidth}
				\begin{itemize}
					\item Avaliar se as estratégias de comunicação estão funcionando.
					\item Interagir com as pessoas para manter o alinhamento.
					\item Determinar as necessidades e o potencial de influência.
					\item Listar pessoas, grupos ou organizações afetadas pelo projeto.
				\end{itemize}
			\end{column}
		\end{columns}
		\begin{enumerate}[a]
			\item A, B, C, D.
			\item \alert<2>{C, D, B, A}
			\item D, C, A, B.
			\item B, A, D, C.
			\item A, C, B, D.
		\end{enumerate}
	\end{enumerate}
\end{frame}

\begin{frame}{Partes Interessadas}
	\begin{enumerate}
		\setcounter{enumi}{3}
		\justifying
		\item Ordene as ações para lidar com uma nova parte interessada influente identificada no meio do projeto:
		\begin{enumerate}[I]
			\item Identificar a nova parte interessada.
			\item Analisar suas expectativas e poder.
			\item Atualizar a estratégia de engajamento.
			\item Executar a interação e monitorar o feedback.
		\end{enumerate}
		\begin{enumerate}[a]
			\item 4, 2, 1, 3.
			\item \alert<2>{1, 2, 3, 4}
			\item 2, 1, 4, 3.
			\item 3, 4, 1, 2.
			\item 1, 3, 2, 4.
		\end{enumerate}
	\end{enumerate}
\end{frame}

\begin{frame}{Partes Interessadas}
	\begin{enumerate}
		\setcounter{enumi}{4}
		\justifying
		\item Durante uma reunião, um stakeholder chave demonstra resistência a uma mudança tecnológica. Para seguir as diretrizes do PMBOK 7, o gerente deve:
		\begin{enumerate}[a]
			\item Ignorar a resistência e seguir o plano original.
			\item Remover o stakeholder das próximas comunicações.
			\item \alert<2>{Escutar ativamente para entender as preocupações e buscar alinhamento.}
			\item Relatar o stakeholder ao RH por comportamento obstrutivo.
			\item Pedir ao patrocinador que force a aceitação da mudança.
		\end{enumerate}
	\end{enumerate}
\end{frame}

\begin{frame}{Partes Interessadas}
	\begin{enumerate}
		\setcounter{enumi}{5}
		\justifying
		\item São categorias de partes interessadas mencionadas no Guia:
		\begin{enumerate}[I]
			\item Clientes e usuários.
			\item Fornecedores e parceiros.
			\item Órgãos reguladores.
			\item A equipe do projeto.
		\end{enumerate}
		\begin{enumerate}[a]
			\item Se apenas I, II e IV estiverem corretas.
			\item Se apenas II e III estiverem corretas.
			\item \alert<2>{Se todas estiverem corretas.}
			\item Se apenas IV estiver correta.
			\item Se apenas I e III estiverem corretas.
		\end{enumerate}
	\end{enumerate}
\end{frame}

%-------------------------------------------------------------------------------
% SUBSEÇÃO 2.2 - DOMÍNIO DE DESEMPENHO DA EQUIPE -------------------------------
%-------------------------------------------------------------------------------
\subsection{Equipe}

\begin{frame}{Equipe}
	\begin{enumerate}
		\justifying
		\item Modelos de desenvolvimento de equipes, como o de Tuckman, descrevem estágios como formação, tormenta, \rule{50pt}{1pt}, desempenho e \rule{50pt}{1pt}.
		\begin{enumerate}[a]
			\item Normatização / Encerramento
			\item \alert<2>{Normalização / Dissolução}
			\item Tormenta / Normatização
			\item Desempenho / Formação
			\item Encerramento / Normatização
		\end{enumerate}
	\end{enumerate}
\end{frame}

\begin{frame}{Equipe}
	\begin{enumerate}
		\setcounter{enumi}{1}
		\justifying
		\item Equipes de alto desempenho demonstram \rule{80pt}{1pt}, o que significa que os membros assumem responsabilidade pelos resultados do projeto, não apenas por suas tarefas individuais.
		\begin{enumerate}[a]
			\item Liderança distribuída
			\item \alert<2>{Responsabilidade Compartilhada}
			\item Autonomia total
			\item Hierarquia clara
			\item Competição saudável
		\end{enumerate}
	\end{enumerate}
\end{frame}

\begin{frame}{Equipe}
	\begin{enumerate}
		\setcounter{enumi}{2}
		\justifying
		\item Ordene as etapas lógicas para resolver um conflito disfuncional na equipe:
		\begin{enumerate}[I]
			\item Identificar a causa raiz do conflito.
			\item Permitir que as partes envolvidas discutam suas perspectivas.
			\item Facilitar a busca por uma solução de ganho mútuo.
			\item Formalizar o acordo e monitorar o relacionamento futuro.
		\end{enumerate}
		\begin{enumerate}[a]
			\item 4, 2, 1, 3.
			\item \alert<2>{1, 2, 3, 4}
			\item 2, 1, 4, 3.
			\item 3, 4, 1, 2.
			\item 1, 3, 2, 4.
		\end{enumerate}
	\end{enumerate}
\end{frame}

\begin{frame}{Equipe}
	\begin{enumerate}
		\setcounter{enumi}{3}
		\justifying
		\item Relacione o conceito de liderança à sua aplicação:
		\begin{columns}
			\begin{column}{0.3\textwidth}
				\begin{enumerate}[A]
					\item Servidora
					\item Transformacional
					\item Situacional
					\item Distribuída
				\end{enumerate}
			\end{column}
			\begin{column}{0.7\textwidth}
				\begin{itemize}
					\item Focar em remover obstáculos para a equipe.
					\item Inspirar através de visão e mudança.
					\item Ajustar o estilo conforme a maturidade da equipe.
					\item Liderança compartilhada entre membros da equipe.
				\end{itemize}
			\end{column}
		\end{columns}
		\begin{enumerate}[a]
			\item A, B, C, D.
			\item \alert<2>{C, A, D, B}
			\item D, C, B, A.
			\item B, D, A, C.
			\item A, C, B, D.
		\end{enumerate}
	\end{enumerate}
\end{frame}

\begin{frame}{Equipe}
	\begin{enumerate}
		\setcounter{enumi}{4}
		\justifying
		\item Quais são traços de uma cultura de equipe positiva?
		\begin{enumerate}[I]
			\item Transparência na comunicação.
			\item Respeito às diferenças individuais.
			\item Aprendizado contínuo com os erros.
			\item Foco exclusivo em métricas individuais.
		\end{enumerate}
		\begin{enumerate}[a]
			\item Se apenas I, II e III estiverem corretas.
			\item Se apenas II e IV estiverem corretas.
			\item \alert<2>{Se apenas I, II e III estiverem corretas.}
			\item Se apenas IV estiver correta.
			\item Se todas estiverem corretas.
		\end{enumerate}
	\end{enumerate}
\end{frame}

\begin{frame}{Equipe}
	\begin{enumerate}
		\setcounter{enumi}{5}
		\justifying
		\item Você percebe que dois desenvolvedores seniores estão em desacordo constante sobre a arquitetura. Seguindo o domínio da Equipe, você deve:
		\begin{enumerate}[a]
			\item Escolher a solução de um deles e impor ao outro.
			\item Ignorar, pois técnicos sempre discutem.
			\item \alert<2>{Facilitar uma discussão técnica baseada em fatos e critérios de sucesso.}
			\item Substituir ambos por profissionais mais dóceis.
			\item Pedir ao cliente para decidir qual arquitetura prefere.
		\end{enumerate}
	\end{enumerate}
\end{frame}

%-------------------------------------------------------------------------------
% SUBSEÇÃO 2.3 - DOMÍNIO DE DESEMPENHO DA ABORDAGEM DE DESENVOLVIMENTO E CICLO DE VIDA
%-------------------------------------------------------------------------------
\subsection{Abordagem de desenvolvimento e ciclo de vida}

\begin{frame}{Abordagem de desenvolvimento e ciclo de vida}
	\begin{enumerate}
		\justifying
		\item As fases do projeto, quando conectadas, formam o \rule{60pt}{1pt}. Este pode ser estruturado de forma sequencial, iterativa ou \rule{60pt}{1pt}.
		\begin{enumerate}[a]
			\item Ciclo de Vida / Linear
			\item \alert<2>{Ciclo de Vida / Sobreposta}
			\item Projeto / Iterativa
			\item Fase / Sequencial
			\item Entrega / Adaptativa
		\end{enumerate}
	\end{enumerate}
\end{frame}

\begin{frame}{Abordagem de desenvolvimento e ciclo de vida}
	\begin{enumerate}
		\setcounter{enumi}{1}
		\justifying
		\item A entrega única ao final do projeto é característica principal da cadência \rule{60pt}{1pt}.
		\begin{enumerate}[a]
			\item Adaptativa
			\item Iterativa
			\item \alert<2>{Preditiva (ou de Entrega Única)}
			\item Híbrida
			\item Incremental
		\end{enumerate}
	\end{enumerate}
\end{frame}

\begin{frame}{Abordagem de desenvolvimento e ciclo de vida}
	\begin{enumerate}
		\setcounter{enumi}{2}
		\justifying
		\item Relacione o termo à sua definição:
		\begin{columns}
			\begin{column}{0.28\textwidth}
				\begin{enumerate}[A]
					\item Cadência
					\item Fase
					\item Iteração
					\item Marco
				\end{enumerate}
			\end{column}
			\begin{column}{0.72\textwidth}
				\begin{itemize}
					\item Ritmo das entregas do projeto.
					\item Grupo de atividades relacionadas que termina com uma entrega.
					\item Ciclo de tempo para desenvolver um incremento de trabalho.
					\item Ponto de decisão ou conclusão significativa.
				\end{itemize}
			\end{column}
		\end{columns}
		\begin{enumerate}[a]
			\item A, B, C, D.
			\item \alert<2>{A, B, C, D}
			\item C, D, A, B.
			\item D, C, B, A.
			\item B, A, D, C.
		\end{enumerate}
	\end{enumerate}
\end{frame}

\begin{frame}{Abordagem de desenvolvimento e ciclo de vida}
	\begin{enumerate}
		\setcounter{enumi}{3}
		\justifying
		\item Ordene as etapas para definir a abordagem de desenvolvimento:
		\begin{enumerate}[I]
			\item Avaliar a natureza do produto e os requisitos.
			\item Analisar as variáveis organizacionais e de mercado.
			\item Selecionar a abordagem (preditiva, híbrida ou adaptativa).
			\item Definir as fases e marcos do ciclo de vida.
		\end{enumerate}
		\begin{enumerate}[a]
			\item 4, 2, 1, 3.
			\item \alert<2>{1, 2, 3, 4}
			\item 2, 1, 4, 3.
			\item 3, 4, 1, 2.
			\item 1, 3, 2, 4.
		\end{enumerate}
	\end{enumerate}
\end{frame}

\begin{frame}{Abordagem de desenvolvimento e ciclo de vida}
	\begin{enumerate}
		\setcounter{enumi}{4}
		\justifying
		\item Um projeto de construção de uma ponte onde o design é fixo e o solo é bem conhecido deve preferencialmente usar:
		\begin{enumerate}[a]
			\item Scrum puro.
			\item Abordagem adaptativa extrema.
			\item \alert<2>{Abordagem preditiva.}
			\item Nenhuma abordagem formal.
			\item Ciclo de vida puramente iterativo sem plano final.
		\end{enumerate}
	\end{enumerate}
\end{frame}

\begin{frame}{Abordagem de desenvolvimento e ciclo de vida}
	\begin{enumerate}
		\setcounter{enumi}{5}
		\justifying
		\item Sobre a abordagem híbrida:
		\begin{enumerate}[I]
			\item Combina métodos ágeis e preditivos.
			\item Pode usar preditivo para hardware e ágil para software.
			\item É útil quando há requisitos estáveis e outros incertos.
			\item É a abordagem mais simples de gerenciar.
		\end{enumerate}
		\begin{enumerate}[a]
			\item Se apenas I, II e III estiverem corretas.
			\item Se apenas II e IV estiverem corretas.
			\item \alert<2>{Se apenas I, II e III estiverem corretas.}
			\item Se apenas IV estiver correta.
			\item Se todas estiverem corretas.
		\end{enumerate}
	\end{enumerate}
\end{frame}

%-------------------------------------------------------------------------------
% SUBSEÇÃO 2.4 - DOMÍNIO DE DESEMPENHO DO PLANEJAMENTO -------------------------
%-------------------------------------------------------------------------------
\subsection{Planejamento}

\begin{frame}{Planejamento}
	\begin{enumerate}
		\justifying
		\item O planejamento de \rule{60pt}{1pt} envolve determinar quais bens e serviços serão adquiridos fora da organização do projeto.
		\begin{enumerate}[a]
			\item Aquisições
			\item \alert<2>{Aquisições}
			\item Recursos
			\item Riscos
			\item Comunicação
		\end{enumerate}
	\end{enumerate}
\end{frame}

\begin{frame}{Planejamento}
	\begin{enumerate}
		\setcounter{enumi}{1}
		\justifying
		\item O planejamento iterativo onde o trabalho próximo é planejado em detalhes e o trabalho futuro em alto nível chama-se planeamento em \rule{60pt}{1pt} \rule{60pt}{1pt}.
		\begin{enumerate}[a]
			\item Ondas Sucessivas
			\item \alert<2>{Ondas Sucessivas}
			\item Iterações Curtas
			\item Planejamento Detalhado
			\item Abordagem Ágil
		\end{enumerate}
	\end{enumerate}
\end{frame}

\begin{frame}{Planejamento}
	\begin{enumerate}
		\setcounter{enumi}{2}
		\justifying
		\item Relacione a técnica de estimativa à sua característica:
		\begin{columns}
			\begin{column}{0.3\textwidth}
				\begin{enumerate}[A]
					\item Análoga
					\item Paramétrica
					\item Bottom-up
					\item Três Pontos
				\end{enumerate}
			\end{column}
			\begin{column}{0.7\textwidth}
				\begin{itemize}
					\item Baseada em projetos similares anteriores.
					\item Usa algoritmos baseados em dados históricos.
					\item Soma estimativas de componentes de nível inferior.
					\item Considera cenários otimista, pessimista e provável.
				\end{itemize}
			\end{column}
		\end{columns}
		\begin{enumerate}[a]
			\item A, B, C, D.
			\item \alert<2>{A, B, C, D}
			\item C, D, A, B.
			\item D, C, B, A.
			\item B, A, D, C.
		\end{enumerate}
	\end{enumerate}
\end{frame}

\begin{frame}{Planejamento}
	\begin{enumerate}
		\setcounter{enumi}{3}
		\justifying
		\item Ordene os passos gerais do planejamento de cronograma:
		\begin{enumerate}[I]
			\item Definir o escopo e as entregas.
			\item Identificar as atividades necessárias.
			\item Sequenciar as atividades e estimar durações.
			\item Desenvolver o cronograma e definir marcos.
		\end{enumerate}
		\begin{enumerate}[a]
			\item 4, 2, 1, 3.
			\item \alert<2>{1, 2, 3, 4}
			\item 2, 1, 4, 3.
			\item 3, 4, 1, 2.
			\item 1, 3, 2, 4.
		\end{enumerate}
	\end{enumerate}
\end{frame}

\begin{frame}{Planejamento}
	\begin{enumerate}
		\setcounter{enumi}{4}
		\justifying
		\item Ao planejar um projeto altamente inovador, o gerente descobre que não há dados históricos. Ele deve priorizar:
		\begin{enumerate}[a]
			\item Estimativas análogas forçadas.
			\item \alert<2>{Planejamento em ondas sucessivas com base em protótipos.}
			\item Ignorar o planejamento e apenas executar.
			\item Usar um cronograma de outro setor totalmente diferente.
			\item Pedir ao patrocinador que defina todas as datas.
		\end{enumerate}
	\end{enumerate}
\end{frame}

%-------------------------------------------------------------------------------
% SUBSEÇÃO 2.5 - DOMÍNIO DE DESEMPENHO DO TRABALHO DO PROJETO ------------------
%-------------------------------------------------------------------------------
\subsection{Trabalho do Projeto}

\begin{frame}{Trabalho do Projeto}
	\begin{enumerate}
		\justifying
		\item A gestão do trabalho do projeto envolve a manutenção de um \rule{70pt}{1pt} equilibrado para evitar sobrecarga ou ociosidade da equipe.
		\begin{enumerate}[a]
			\item Fluxo de Trabalho
			\item \alert<2>{Fluxo de Trabalho}
			\item Cronograma Detalhado
			\item Orçamento Balanceado
			\item Capacidade Máxima
		\end{enumerate}
	\end{enumerate}
\end{frame}

\begin{frame}{Trabalho do Projeto}
	\begin{enumerate}
		\setcounter{enumi}{1}
		\justifying
		\item O ato de compartilhar conhecimento e insights durante a execução do projeto para melhorar processos futuros é chamado de \rule{80pt}{1pt}.
		\begin{enumerate}[a]
			\item Aprendizado Contínuo (ou Gestão do Conhecimento)
			\item \alert<2>{Aprendizado Contínuo (ou Gestão do Conhecimento)}
			\item Retrospectiva Final
			\item Relatório de Lições Aprendidas
			\item Transferência de Conhecimento
		\end{enumerate}
	\end{enumerate}
\end{frame}

\begin{frame}{Trabalho do Projeto}
	\begin{enumerate}
		\setcounter{enumi}{2}
		\justifying
		\item Relacione os conceitos de gestão de trabalho:
		\begin{columns}
			\begin{column}{0.3\textwidth}
				\begin{enumerate}[A]
					\item Gargalo
					\item Lead Time
					\item Throughput
					\item WIP
				\end{enumerate}
			\end{column}
			\begin{column}{0.7\textwidth}
				\begin{itemize}
					\item Ponto de obstrução que limita o fluxo de saída.
					\item Tempo total do início ao fim de uma tarefa.
					\item Quantidade de itens entregues em um período.
					\item Quantidade de trabalho em progresso.
				\end{itemize}
			\end{column}
		\end{columns}
		\begin{enumerate}[a]
			\item A, B, C, D.
			\item \alert<2>{A, B, C, D}
			\item C, D, A, B.
			\item D, C, B, A.
			\item B, A, D, C.
		\end{enumerate}
	\end{enumerate}
\end{frame}

\begin{frame}{Trabalho do Projeto}
	\begin{enumerate}
		\setcounter{enumi}{3}
		\justifying
		\item Ordene as ações de gerenciamento de aquisições (parte do trabalho do projeto):
		\begin{enumerate}[I]
			\item Definir a necessidade e os termos de referência.
			\item Selecionar o fornecedor e assinar o contrato.
			\item Gerenciar o desempenho do fornecedor durante o trabalho.
			\item Encerrar o contrato após a entrega e aceitação.
		\end{enumerate}
		\begin{enumerate}[a]
			\item 4, 2, 1, 3.
			\item \alert<2>{1, 2, 3, 4}
			\item 2, 1, 4, 3.
			\item 3, 4, 1, 2.
			\item 1, 3, 2, 4.
		\end{enumerate}
	\end{enumerate}
\end{frame}

\begin{frame}{Trabalho do Projeto}
	\begin{enumerate}
		\setcounter{enumi}{4}
		\justifying
		\item A equipe relata que uma ferramenta de software está travando e atrasando as tarefas. No domínio do Trabalho do Projeto, o gerente deve:
		\begin{enumerate}[a]
			\item Dizer à equipe para trabalhar no final de semana para compensar.
			\item Ignorar, pois falhas de software são normais.
			\item \alert<2>{Atuar para resolver o problema técnico ou substituir a ferramenta visando a eficiência do fluxo.}
			\item Pedir ao cliente para estender o prazo sem explicar o motivo.
			\item Responsabilizar o membro que descobriu o erro.
		\end{enumerate}
	\end{enumerate}
\end{frame}

\begin{frame}{Trabalho do Projeto}
	\begin{enumerate}
		\setcounter{enumi}{5}
		\justifying
		\item Sobre o gerenciamento de recursos físicos:
		\begin{enumerate}[I]
			\item Envolve logística de entrega de materiais.
			\item Requer controle de desperdícios.
			\item Não se aplica a projetos de software.
			\item Inclui a segurança do local de trabalho.
		\end{enumerate}
		\begin{enumerate}[a]
			\item Se apenas I, II e IV estiverem corretas.
			\item Se apenas II e III estiverem corretas.
			\item \alert<2>{Se apenas I, II e IV estiverem corretas.}
			\item Se apenas IV estiver correta.
			\item Se todas estiverem corretas.
		\end{enumerate}
	\end{enumerate}
\end{frame}

%-------------------------------------------------------------------------------
% SUBSEÇÃO 2.6 - DOMÍNIO DE DESEMPENHO DA ENTREGA ------------------------------
%-------------------------------------------------------------------------------
\subsection{Entrega}

\begin{frame}{Entrega}
	\begin{enumerate}
		\justifying
		\item O diferencial entre 'Entregas' e 'Resultados' é:
		\begin{enumerate}[a]
			\item Não há diferença; são sinônimos.
			\item \alert<2>{Entregas são produtos tangíveis; Resultados são os benefícios e efeitos gerados por essas entregas.}
			\item Resultados são documentos; Entregas são máquinas.
			\item Entregas são feitas pela equipe; Resultados são feitos pelo cliente.
			\item Resultados só ocorrem em projetos ágeis.
		\end{enumerate}
	\end{enumerate}
\end{frame}

\begin{frame}{Entrega}
	\begin{enumerate}
		\setcounter{enumi}{1}
		\justifying
		\item A conformidade com os \rule{50pt}{1pt} de \rule{50pt}{1pt} garante que a entrega atende às necessidades das partes interessadas para ser formalmente aprovada.
		\begin{enumerate}[a]
			\item Critérios / Aceitação
			\item \alert<2>{Critérios / Aceitação}
			\item Requisitos / Escopo
			\item Especificações / Qualidade
			\item Padrões / Entrega
		\end{enumerate}
	\end{enumerate}
\end{frame}

\begin{frame}{Entrega}
	\begin{enumerate}
		\setcounter{enumi}{2}
		\justifying
		\item Relacione os termos de entrega:
		\begin{columns}
			\begin{column}{0.28\textwidth}
				\begin{enumerate}[A]
					\item Elicitação
					\item Validação
					\item Verificação
					\item Backlog
				\end{enumerate}
			\end{column}
			\begin{column}{0.72\textwidth}
				\begin{itemize}
					\item Coletar requisitos com os stakeholders.
					\item Garantir que o produto atende ao uso pretendido.
					\item Garantir que o produto atende às especificações técnicas.
					\item Lista priorizada de trabalhos a serem entregues.
				\end{itemize}
			\end{column}
		\end{columns}
		\begin{enumerate}[a]
			\item A, B, C, D.
			\item \alert<2>{A, B, D, C}
			\item C, D, A, B.
			\item D, C, B, A.
			\item B, A, D, C.
		\end{enumerate}
	\end{enumerate}
\end{frame}

\begin{frame}{Entrega}
	\begin{enumerate}
		\setcounter{enumi}{3}
		\justifying
		\item Ordene o fluxo de uma entrega típica:
		\begin{enumerate}[I]
			\item Elicitação e análise de requisitos.
			\item Desenvolvimento/Construção da entrega.
			\item Verificação da qualidade interna.
			\item Validação e aceitação pelo cliente.
		\end{enumerate}
		\begin{enumerate}[a]
			\item 4, 2, 1, 3.
			\item \alert<2>{1, 2, 3, 4}
			\item 2, 1, 4, 3.
			\item 3, 4, 1, 2.
			\item 1, 3, 2, 4.
		\end{enumerate}
	\end{enumerate}
\end{frame}

\begin{frame}{Entrega}
	\begin{enumerate}
		\setcounter{enumi}{4}
		\justifying
		\item O custo associado a evitar falhas (como treinamento e processos robustos) é chamado de Custo de \rule{60pt}{1pt}.
		\begin{enumerate}[a]
			\item Avaliação
			\item \alert<2>{Prevenção}
			\item Falha
			\item Inspeção
			\item Correção
		\end{enumerate}
	\end{enumerate}
\end{frame}

\begin{frame}{Entrega}
	\begin{enumerate}
		\setcounter{enumi}{5}
		\justifying
		\item Um cliente solicita uma nova funcionalidade que não estava no plano original, mas que agregaria muito valor. Em um ambiente adaptativo, o gerente deve:
		\begin{enumerate}[a]
			\item Negar, pois o plano já foi assinado.
			\item \alert<2>{Adicionar ao backlog e priorizar junto ao cliente para a próxima iteração.}
			\item Fazer imediatamente, parando todo o resto sem avisar ninguém.
			\item Cobrar uma multa punitiva pelo pedido.
			\item Esperar o final do projeto para discutir isso.
		\end{enumerate}
	\end{enumerate}
\end{frame}

\begin{frame}{Entrega}
	\begin{enumerate}
		\setcounter{enumi}{6}
		\justifying
		\item Fatores que indicam uma entrega bem-sucedida:
		\begin{enumerate}[I]
			\item Satisfação do cliente.
			\item Conformidade com os requisitos.
			\item Alcance dos benefícios de negócio.
			\item Ausência de qualquer mudança durante a execução.
		\end{enumerate}
		\begin{enumerate}[a]
			\item Se apenas I, II e III estiverem corretas.
			\item Se apenas II e IV estiverem corretas.
			\item \alert<2>{Se apenas I, II e III estiverem corretas.}
			\item Se apenas IV estiver correta.
			\item Se todas estiverem corretas.
		\end{enumerate}
	\end{enumerate}
\end{frame}

%-------------------------------------------------------------------------------
% SUBSEÇÃO 2.7 - DOMÍNIO DE DESEMPENHO DA MEDIÇÃO ------------------------------
%-------------------------------------------------------------------------------
\subsection{Medição}

\begin{frame}{Medição}
	\begin{enumerate}
		\justifying
		\item No Gerenciamento de Valor Agregado (EVA), o \rule{70pt}{1pt} (CPI) mede a eficiência de custo do projeto.
		\begin{enumerate}[a]
			\item Índice de Desempenho de Custos
			\item \alert<2>{Índice de Desempenho de Custos}
			\item Índice de Desempenho do Cronograma
			\item Valor Planejado
			\item Variância de Custo
		\end{enumerate}
	\end{enumerate}
\end{frame}

\begin{frame}{Medição}
	\begin{enumerate}
		\setcounter{enumi}{1}
		\justifying
		\item A diferença entre o desempenho planejado e o desempenho real é chamada de \rule{60pt}{1pt}.
		\begin{enumerate}[a]
			\item Variação
			\item \alert<2>{Variação}
			\item Desvio
			\item Correção
			\item Baseline
		\end{enumerate}
	\end{enumerate}
\end{frame}

\begin{frame}{Medição}
	\begin{enumerate}
		\setcounter{enumi}{2}
		\justifying
		\item Relacione a métrica ao seu foco:
		\begin{columns}
			\begin{column}{0.3\textwidth}
				\begin{enumerate}[A]
					\item SPI
					\item Burn-down Chart
					\item Ciclo de Vida do Defeito
					\item NPS
				\end{enumerate}
			\end{column}
			\begin{column}{0.7\textwidth}
				\begin{itemize}
					\item Eficiência do cronograma.
					\item Progresso do trabalho restante em ágil.
					\item Qualidade do produto.
					\item Satisfação do cliente.
				\end{itemize}
			\end{column}
		\end{columns}
		\begin{enumerate}[a]
			\item A, B, C, D.
			\item \alert<2>{A, B, C, D}
			\item C, D, A, B.
			\item D, C, B, A.
			\item B, A, D, C.
		\end{enumerate}
	\end{enumerate}
\end{frame}

\begin{frame}{Medição}
	\begin{enumerate}
		\setcounter{enumi}{3}
		\justifying
		\item Ordene as fases do ciclo de medição:
		\begin{enumerate}[I]
			\item Definir o que medir e estabelecer a linha de base.
			\item Coletar os dados reais.
			\item Analisar as variações e tendências.
			\item Implementar ações corretivas ou preventivas.
		\end{enumerate}
		\begin{enumerate}[a]
			\item 4, 2, 1, 3.
			\item \alert<2>{1, 2, 3, 4}
			\item 2, 1, 4, 3.
			\item 3, 4, 1, 2.
			\item 1, 3, 2, 4.
		\end{enumerate}
	\end{enumerate}
\end{frame}

\begin{frame}{Medição}
	\begin{enumerate}
		\setcounter{enumi}{4}
		\justifying
		\item Um projeto apresenta CPI de 0,85 e SPI de 1,10. Isso significa que o projeto está:
		\begin{enumerate}[a]
			\item Abaixo do orçamento e adiantado.
			\item Acima do orçamento e atrasado.
			\item \alert<2>{Acima do orçamento e adiantado.}
			\item Abaixo do orçamento e atrasado.
			\item Dentro do planejado em ambos.
		\end{enumerate}
	\end{enumerate}
\end{frame}

\begin{frame}{Medição}
	\begin{enumerate}
		\setcounter{enumi}{5}
		\justifying
		\item Sobre a apresentação de dados (Reporting):
		\begin{enumerate}[I]
			\item Deve ser adaptada ao público-alvo.
			\item Deve focar no que é acionável.
			\item Deve ser feita diariamente para todos os stakeholders.
			\item Deve utilizar recursos visuais sempre que possível.
		\end{enumerate}
		\begin{enumerate}[a]
			\item Se apenas I, II e IV estiverem corretas.
			\item Se apenas II e III estiverem corretas.
			\item \alert<2>{Se apenas I, II e IV estiverem corretas.}
			\item Se apenas IV estiver correta.
			\item Se todas estiverem corretas.
		\end{enumerate}
	\end{enumerate}
\end{frame}

%-------------------------------------------------------------------------------
% SUBSEÇÃO 2.8 - DOMÍNIO DE DESEMPENHO DA INCERTEZA ----------------------------
%-------------------------------------------------------------------------------
\subsection{Incerteza}

\begin{frame}{Incerteza}
	\begin{enumerate}
		\justifying
		\item A resposta a um risco negativo que visa reduzir a probabilidade ou o impacto é chamada de \rule{60pt}{1pt}.
		\begin{enumerate}[a]
			\item Mitigação
			\item \alert<2>{Mitigação}
			\item Transferência
			\item Evitação
			\item Aceitação
		\end{enumerate}
	\end{enumerate}
\end{frame}

\begin{frame}{Incerteza}
	\begin{enumerate}
		\setcounter{enumi}{1}
		\justifying
		\item Quais são estratégias de resposta a oportunidades (riscos positivos)?
		\begin{enumerate}[I]
			\item Explorar.
			\item Compartilhar.
			\item Melhorar (Enhance).
			\item Evitar.
		\end{enumerate}
		\begin{enumerate}[a]
			\item Se apenas I, II e III estiverem corretas.
			\item Se apenas II e IV estiverem corretas.
			\item \alert<2>{Se apenas I, II e III estiverem corretas.}
			\item Se apenas IV estiver correta.
			\item Se todas estiverem corretas.
		\end{enumerate}
	\end{enumerate}
\end{frame}

\begin{frame}{Incerteza}
	\begin{enumerate}
		\setcounter{enumi}{2}
		\justifying
		\item Relacione o termo à sua característica:
		\begin{columns}
			\begin{column}{0.32\textwidth}
				\begin{enumerate}[A]
					\item Apetite ao Risco
					\item Limite de Risco
					\item Reserva de Contingência
					\item Risco Residual
				\end{enumerate}
			\end{column}
			\begin{column}{0.68\textwidth}
				\begin{itemize}
					\item Grau de incerteza que a organização está disposta a aceitar.
					\item Valor específico que a organização aceita arriscar.
					\item Tempo ou dinheiro alocado para riscos identificados.
					\item Risco que permanece após a implementação das respostas.
				\end{itemize}
			\end{column}
		\end{columns}
		\begin{enumerate}[a]
			\item A, B, C, D.
			\item \alert<2>{A, B, C, D}
			\item C, D, A, B.
			\item D, C, B, A.
			\item B, A, D, C.
		\end{enumerate}
	\end{enumerate}
\end{frame}

\begin{frame}{Incerteza}
	\begin{enumerate}
		\setcounter{enumi}{3}
		\justifying
		\item Ordene as etapas da gestão de riscos:
		\begin{enumerate}[I]
			\item Identificar os riscos potenciais.
			\item Realizar análise qualitativa e quantitativa.
			\item Planejar e implementar as respostas aos riscos.
			\item Monitorar e revisar o status dos riscos continuamente.
		\end{enumerate}
		\begin{enumerate}[a]
			\item 4, 2, 1, 3.
			\item \alert<2>{1, 2, 3, 4}
			\item 2, 1, 4, 3.
			\item 3, 4, 1, 2.
			\item 1, 3, 2, 4.
		\end{enumerate}
	\end{enumerate}
\end{frame}

\begin{frame}{Incerteza}
	\begin{enumerate}
		\setcounter{enumi}{4}
		\justifying
		\item Quando um risco ocorre, ele deixa de ser uma incerteza e passa a ser considerado um \rule{60pt}{1pt}.
		\begin{enumerate}[a]
			\item Problema (ou Questão/Issue)
			\item \alert<2>{Problema (ou Questão/Issue)}
			\item Falha
			\item Mudança
			\item Incidente
		\end{enumerate}
	\end{enumerate}
\end{frame}

\begin{frame}{Incerteza}
	\begin{enumerate}
		\setcounter{enumi}{5}
		\justifying
		\item Um furacão destrói o canteiro de obras. Este evento não estava mapeado individualmente, mas a empresa tinha um seguro. A estratégia usada foi:
		\begin{enumerate}[a]
			\item Evitar.
			\item Mitigar.
			\item \alert<2>{Transferir.}
			\item Explorar.
			\item Aceitar passivamente.
		\end{enumerate}
	\end{enumerate}
\end{frame}

\begin{frame}{Incerteza}
	\begin{enumerate}
		\setcounter{enumi}{6}
		\justifying
		\item Fontes de incerteza em projetos incluem:
		\begin{enumerate}[I]
			\item Variabilidade de processos.
			\item Ambiguidade de requisitos.
			\item Mudanças de stakeholders.
			\item Complexidade de integrações.
		\end{enumerate}
		\begin{enumerate}[a]
			\item Se todas estiverem corretas.
			\item Se apenas I e II estiverem corretas.
			\item \alert<2>{Se todas estiverem corretas.}
			\item Se apenas I, II e III estiverem corretas.
			\item Se apenas II e III estiverem corretas.
		\end{enumerate}
	\end{enumerate}
\end{frame}

%===============================================================================
% FIM ==========================================================================
%===============================================================================

\begin{frame}
	\centering
	\Large
	FIM
\end{frame}

\end{document}